\section{Wnioski}
Rozdział wnioski nie może składać się z podrozdziałów; powinien składać się z takich
samych akapitów jak Wstęp, ale napisanych w czasie przeszłym i inaczej sformułowanych:

I akapit: w zależności od tematyki pracy: ogólny stan rynku.

II akapit: główne cele pracy (podkreślić, że zostały zrealizowane) .

III akapit: nowe ujęcie w zakresie pracy: Wyszczególnienie nowego ujęcia w omawianym
akapicie można przyjąć postać: "Do pracy wprowadzono następujące nowe
rozwiązania..." lub "Elementem nowości wprowadzonym do pracy jest..." lub „Nowym
ujęciem tematyki poruszanej w pracy jest …”

IV akapit: podczas realizacji pracy, autor napotkał na kilka problemów; wymienić dwa,
trzy i w kilku zdaniach napisać jak każdy z nich został rozwiązany

V akapit: "Praca napisana jest w formie pozwalającej na jej rozbudowę." Tutaj trzeba
dodać, jakie elementy czy moduły mogłyby (według autora) uatrakcyjnić pracę.
Wymienić dwa, trzy i krótko opisać.
